
\begin{corolario}[Dos casillas]\label{cor:dos-casillas}
Sea $M$ una máscara que cumple \textup{(H)} con constantes $(C_p,\gamma_A)$, sean
$k_1'=k_1(\tfrac54C_p,\tfrac34\gamma_A)$ y $k\ge k_2:=\max\{2k_1,k_1'+1\}$. Sean $e,f$ casillas en filas
y columnas distintas, $\pi^T=w_ew_fe^{L_T(M_{ef})-L_T(M)}$ y
$\pi_1=p_ep_f\{1-\mathsf G(e,f)+\hat c_e+\hat c_f\}$. Entonces
\[
 |\pi^T-\pi_1|\le C_2'\,\frac{p_ep_f}{k^2}\le\frac{C_2}{k^4},
\]
con $C_2',C_2$ que sólo dependen de $(C_p,\gamma_A)$. Con más precisión, con primas para los datos de
$M_e$,
\begin{equation}\label{eq:producto-inserciones}
 \pi^T=p_e(1+\hat c_e)e^{r_e(M)}\,p'_f(1+\hat c'_f)e^{r_f(M_e)},\qquad
 \Bigl|\frac{p'_f}{p_f}-1+\mathsf G(e,f)\Bigr|\le\frac C{k^2},\qquad
 |\hat c'_f-\hat c_f|\le\frac C{k^2},
\end{equation}
y $|G'_{ab}-G^0_{ab}|\le C/k^2$ para las líneas $a,b$ de $M_e$, donde $G'$ es la matriz de Green de
$M_e$. En la clase $\mathcal K$ estas cotas valen con $(C_p,\gamma_A)=(\tfrac32M_0,0{,}88\gamma)$ y $k$
por encima de $k_{\mathcal K}$ y de $k_2$.
\end{corolario}
\begin{proof}
Usamos la notación de la demostración del Teorema~\ref{teo:insercion-torres} con $e_0=e$; las constantes
$C$ sólo dependen de $(C_p,\gamma_A)$.

\emph{La matriz de Green del hijo.} En $v_I^\perp$ se tiene $S'(\tau')^+=(\mathsf S-W)^+=G'-\Pi'$ y la
identidad de la resolvente $(\mathsf S-W)^+-\mathsf S^+=\mathsf S^+W(\mathsf S-W)^+$.
Por~\eqref{eq:green-torres} para el hijo, que cumple \textup{(H)} con
$(\tfrac54C_p,\tfrac34\gamma_A)$, las entradas de $S'(\tau')^+$ difieren de $\delta_{db}$ en
$O(k^{-1})$, y $\|\mathsf S^+\|_\infty\le4K_G$. Por el Paso~7 de aquella demostración,
$\max|W_{ab}|=O(k^{-2})$ y $\|W\|_\infty=O(k^{-1})$. Además $G^0_{II}v_I=v_I-q$ con $q=G^0_{IJ}v_J$,
$|q_b|\le2g_G/k$ y $\sum_b|q_b|\le2K_G$, de modo que
\[
 \mathsf S^++\Pi'-G^0_{II}=\frac{v_Iq^\top+qv_I^\top}{|v_I|^2}-\frac{(v_I^\top q)\,v_Iv_I^\top}{|v_I|^4}
\]
tiene entradas $O(k^{-2})$. En conjunto resulta $|G'_{ab}-G^0_{ab}|\le C/k^2$ para $a,b\in I$.

\emph{La marginal del hijo.} Como $a_f^\top\zeta=s_f+\mathsf G(e,f)$,
\[
 p'_f-p_f\bigl(1-\mathsf G(e,f)\bigr)=-\varsigma_f\mathsf G(e,f)+\tau_fa_f^\top\zeta
   +(\Delta p_f-\tau_fs_f).
\]
Los tres términos son $O(p_fk^{-2})$ por $\varsigma_f\le\tfrac43p_f^2$, $|\mathsf G(e,f)|\le4g_G/k$,
\eqref{eq:cotas-hijo} y $|\Delta p_f-\tau_fs_f|\le0{,}9\tau_fs_f^2$. Es la segunda cota
de~\eqref{eq:producto-inserciones}.

\emph{La corrección del hijo.} Por~\eqref{eq:chat-torres}, $\hat c_f=\tau_fX_f/(2p_f)$. La diferencia
$X'_f-X_f$ tiene tres partes. Las casillas de $F$ faltan en la suma de $M_e$; dos de ellas comparten
línea con $f$, y en las demás $|\mathsf G(f,e'')|\le4g_G/k$, de modo que con
$|\varepsilon|\le(4g_G+5C_p)/k$ y las masas de línea $\le\tfrac76$ esta parte es $O(k^{-2})$. La
variación de $\varepsilon_{e''}$ en las casillas comunes es $O(k^{-2})$, porque $\ddot c_2$ y
$\ddot\tau$ están acotadas en el segmento, $\Delta p=O(pk^{-1})$ y
$|\mathsf G'(e'',e'')-\mathsf G(e'',e'')|\le4C/k^2$. La variación de los productos
$\tau_{e''}\varepsilon_{e''}\mathsf G(f,e'')$ se controla con $|\tau'-\tau|\le1{,}8\tau\cdot6K_Ga$ y con
$\sum_{e''}\tau_{e''}|\mathsf G(f,e'')|\le\tfrac73K_G$. Por último,
$|\Delta(\tau_f/p_f)|\le1{,}1|\Delta p_f|$. Por tanto $|\hat c'_f-\hat c_f|\le C/k^2$.

\emph{Conclusión.} Como $(M_e)_f=M_{ef}$, el cociente $\pi^T$ es el producto de los dos cocientes de
inserción $w_ee^{L_T(M_e)-L_T(M)}$ y $w_fe^{L_T(M_{ef})-L_T(M_e)}$. El
Teorema~\ref{teo:insercion-torres} se aplica a $M$ y, por el Paso~4 de su demostración, a $M_e$ con
constantes $(\tfrac54C_p,\tfrac34\gamma_A)$ en tamaño $k-1\ge k_1'$. Esto da el producto
de~\eqref{eq:producto-inserciones}. Escribiendo
$p'_f(1+\hat c'_f)=p_f(1-\mathsf G(e,f)+\theta_1)(1+\hat c_f+\theta_2)$ con
$|\theta_1|,|\theta_2|\le C/k^2$, obtenemos $\pi^T-\pi_1=p_ep_f\Psi$ con
\begin{equation}\label{eq:psi-dos-casillas}
 \Psi=(1+\hat c_e)\bigl(1-\mathsf G(e,f)+\theta_1\bigr)(1+\hat c_f+\theta_2)e^{r_e(M)+r_f(M_e)}
      -\bigl(1-\mathsf G(e,f)+\hat c_e+\hat c_f\bigr).
\end{equation}
Los términos de $\Psi$ lineales en $(\hat c_e,\hat c_f,\mathsf G(e,f))$ se cancelan. Los productos
$\mathsf G\hat c$ y $\hat c\hat c$ son $O(k^{-2})$, y también lo son $\theta_1,\theta_2,r_e(M),r_f(M_e)$.
Así $|\Psi|\le C_2'/k^2$, y $p_ep_f\le C_p^2/k^2$ da la segunda desigualdad.
\end{proof}
